\chapter{Rewrite Strategies and Efficient Rule Engines}
\label{ch:rewrite-strategies}

A correct rule collection is not yet an optimizer.
The implementation must find applicable rules, select an occurrence, decide which facts justify side conditions, and schedule the next search after a change.
A full-tree rescan is easy to understand but can spend most of its time rediscovering that unrelated nodes have not changed.
Efficient engines make those dependencies explicit without weakening the rule contract.

\section{Make strategy a first-class policy}
\label{sec:ch85-strategy}
An outermost strategy prefers positions nearest the root; an innermost strategy prefers positions without reducible descendants.
Leftmost order breaks ties between incomparable positions.
Rule priority breaks ties between patterns matching the same occurrence.
These choices affect trace length, intermediate size, and sometimes whether a normal form is reached at all.
They are not profitability proofs.

The companion visits positions in preorder for its default strategy and sorts positions by decreasing depth for its alternate innermost-first policy.
The latter is deepest-first with stable traversal ties; it should not be misdescribed as every possible conventional leftmost-innermost schedule.
Table~\ref{tab:ch85-strategies} states what each policy controls.
A budget is independent of these choices and can stop either policy before completion.

\begin{table}[htbp]
\caption{A deterministic engine specifies position, rule, and stopping policies separately.}
\label{tab:ch85-strategies}
\centering
\begin{tabular}{ll}
\toprule
Policy & Question answered \\
\midrule
position order & which redex location is considered first? \\
rule priority & which matching rule wins at that location? \\
side-condition order & which analyses are consulted and when? \\
revisit policy & which nodes become candidates after a change? \\
stopping policy & fixed point, phase boundary, or budget exhaustion? \\
\bottomrule
\end{tabular}
\end{table}

For the pure identity system, both executed strategies reach v from any finite nesting of add-zero and multiply-one wrappers.
The test checks exact trace length and replay for 50 nesting depths.
That evidence does not establish that every added rule or effectful operation remains strategy-independent.
An engine should attach its strategy version to reproducibility records when selected output is not canonical.

\section{Index patterns without skipping candidates}
\label{sec:ch85-index}
A root-symbol/arity index is a simple first filter for nonvariable left sides.
Rules rooted at mul cannot match an add node, so they need not enter the full matcher.
Deeper discrimination trees encode successive symbol tests across patterns and share common tests.
Decision DAGs additionally share repeated tests or suffix decisions when their binding environments agree.

Figure~\ref{fig:ch85-dispatch} separates cheap structural discrimination from variable consistency and semantic conditions.
A structural filter can admit false candidates but must not omit any actual match.
For nonlinear patterns such as pair(x,x), testing both child shapes still does not establish equality of the two matched subterms.
The full matcher must retain the repeated-variable constraint.

\begin{figure}[htbp]
\centering
\begin{tikzpicture}[node distance=12mm]
\node[basic block] (root) {root symbol and arity index};
\node[basic block,below=of root] (shape) {shared structural tests};
\node[basic block,below=of shape] (bind) {consistent metavariable bindings};
\node[basic block,below=of bind] (side) {semantic side conditions};
\draw[cfg edge] (root)--(shape);\draw[cfg edge] (shape)--(bind);\draw[cfg edge] (bind)--(side);
\end{tikzpicture}
\caption{Pattern indexing narrows structural candidates but does not discharge semantic premises.}
\label{fig:ch85-dispatch}
\end{figure}

For $N$ nodes and $R$ rules, naive discovery attempts up to $NR$ matches per scan.
Root indexing replaces $R$ with the candidate bucket size at each node, but a worst-case bucket can still contain every rule.
More elaborate indexes have construction and update costs that should be amortized over the number of queries.
The lab intentionally executes the naive matcher so its candidate completeness is inspectable; the chapter's indexed design is an engineering extension, not a benchmark claim about the reference implementation.

Algorithm~\ref{alg:ch85-index} sketches safe indexed discovery.
Its cost is index lookup plus the sum of attempted match and side-condition costs.
Expensive conditions such as dominance or alias queries should run after cheap shape tests when this does not alter required diagnostic behavior.
Caching their answers requires the analysis epoch in the cache key.

\begin{algorithm}[htbp]
\caption{Discover candidates through an overapproximating pattern index.}
\label{alg:ch85-index}
\begin{pseudocode}
\Require node, complete pattern index, current analysis epoch
\Ensure every admitted rewrite has a consistent match and valid premises
\State retrieve candidate rules by root and structural discrimination
\ForAll{candidates in declared priority order}
\State run complete matcher, including repeated-variable checks
\If{matching succeeds}
\State query all semantic conditions against the current IR and epoch
\If{every condition is established} \State yield rule and substitution \EndIf
\EndIf
\EndFor
\end{pseudocode}
\end{algorithm}

\section{Revisit the dependency neighborhood}
\label{sec:ch85-worklist}
Replacing a node can expose a new parent pattern, invalidate a use-count condition, or remove the last reference to a child.
An affected-node worklist is sound only if its dependency relation covers all such changes.
For immutable trees, rebuilding ancestors identifies the direct structural neighborhood.
For shared graphs, reverse-use lists identify every parent, not just the path through which the rewrite was found.

Figure~\ref{fig:ch85-users} shows a shared node whose two parents both need reconsideration.
Failing to enqueue one can miss an optimization; using a stale side-condition fact can do worse and authorize an invalid rewrite.
Worklist duplicate suppression is safe when an already-enqueued node will observe the latest state at processing time.
If the queue instead stores snapshots, it must detect and discard stale snapshots or revalidate them before commit.

\begin{figure}[htbp]
\centering
\begin{tikzpicture}[node distance=14mm and 16mm]
\node[basic block] (n) {changed shared node};
\node[basic block,above left=of n] (p) {parent p};\node[basic block,above right=of n] (q) {parent q};
\draw[cfg edge] (p)--(n);\draw[cfg edge] (q)--(n);
\end{tikzpicture}
\caption{A shared update invalidates every dependent parent, not just one traversal path.}
\label{fig:ch85-users}
\end{figure}

Algorithm~\ref{alg:ch85-worklist} makes validation part of commit.
For $P$ processed candidates, total cost is the sum of discovery, condition, rewrite, and invalidation costs over those candidates.
There is no general linear bound in original graph size when rewrites add nodes or reenable one another.
If a known potential decreases, it can bound successful changes, but failed candidate checks still require accounting.

\begin{algorithm}[htbp]
\caption{Run an affected-node rewrite worklist with stale-match protection.}
\label{alg:ch85-worklist}
\begin{pseudocode}
\Require valid IR, complete dependency neighborhood policy, resource budget
\Ensure valid rewritten IR and explicit fixed-point or incomplete status
\State enqueue initial candidates with version information
\While{worklist nonempty and budget remains}
\State take candidate; revalidate its node and required analyses
\State discover an admitted rewrite under the selected strategy
\If{a rewrite is selected}
\State construct replacement; check invariants; commit change
\State invalidate dependent facts and enqueue every affected neighborhood
\EndIf
\EndWhile
\State distinguish empty-worklist completion from resource exhaustion
\end{pseudocode}
\end{algorithm}

\section{Control growth and phase interaction}
\label{sec:ch85-growth}
Inlining, distributivity, and specialization can increase representation size.
A bounded engine may cap node count, estimated code size, or accumulated cost, but must reject an expansion before losing its last valid representation.
Persistent originals and transactional replacement support rollback when a candidate exceeds the budget or fails verification.
An incomplete pass can still produce a correct improved program if every committed step was sound.

Table~\ref{tab:ch85-phases} contrasts destructive and alternative-preserving designs.
Equality saturation avoids committing immediately to one equivalent form but can consume large memory while retaining alternatives.
Its budget does not prove global optimality of extraction, and tree-cost extraction does not automatically minimize shared machine-code cost.
Combining the approaches can be useful: destructive canonicalization for cheap identities, bounded alternative search for a difficult region, then target-aware selection.

\begin{table}[htbp]
\caption{Search organization moves costs without replacing semantic validation.}
\label{tab:ch85-phases}
\centering
\begin{tabular}{p{0.21\linewidth}p{0.29\linewidth}p{0.35\linewidth}}
\toprule
Organization & Retains & Main risk \\
\midrule
greedy rewriting & selected current form & phase-order opportunity loss \\
phased rewriting & selected form per phase & oscillation across repeated phases \\
backtracking search & explicit alternatives & exponential search growth \\
equality saturation & compact equivalence alternatives & graph growth and extraction cost \\
\bottomrule
\end{tabular}
\end{table}

\section{Diagnose the rule, premise, and representation separately}
\label{sec:ch85-diagnostics}
A failed rewrite check should distinguish no structural match, unmet semantic premise, stale analysis, invariant violation, and resource refusal.
These outcomes have different remedies.
Counting all of them as a generic pattern failure hides whether the rule was never relevant or relevant but unsafe under current facts.
A trace should also record the arithmetic model and relevant flags, not only a printed before/after expression.

Production compiler rule descriptions may generate matchers and combine rules with handwritten lowering code.
The generated code still inherits the semantic assumptions of the rule definitions and the correctness of its index generator.
Testing every small input for a bounded arithmetic domain gives strong finite evidence for that domain; it does not justify replacing the domain with floating point, checked integers, or effectful expressions.
Chapter~\ref{ch:rewrite-optimization} makes this validation boundary concrete.

\section{Preserve conditional facts across a successful rewrite}
\label{sec:ch85-premises}
Suppose a rule replaces a checked arithmetic operation with an unchecked one when a range analysis proves that overflow is impossible.
The proof is about particular operands at a particular program point under a particular analysis state.
It is not a permanent property of the printed opcode.
If an earlier rewrite changes an operand or control predecessor, the cached range premise may no longer apply even though the node ID remains stable.
A candidate record should therefore identify the facts it consumed and the version under which they were established.

A conservative engine invalidates all affected facts after any edit that their preservation policy does not cover.
An incremental engine can update them, but that update procedure becomes part of the soundness argument.
For must-facts such as guaranteed nonnull or guaranteed range bounds, a merge must weaken facts appropriately; reusing the union operation from a may-analysis can accidentally strengthen a premise.
The distinction between structural cache hits and semantic validity remains essential even when both are stored in hash tables.

Conditional rules can also depend on the absence of uses.
Replacing a node changes the use counts of its old and new operands, potentially enabling or disabling rules elsewhere.
A worklist that only revisits the replacement's immediate parents is incomplete for this dependency.
Production engines often categorize rule dependencies so a use-count change can trigger a smaller but complete neighborhood than a whole-function rescan.
The lab's full rescan avoids this indexing subtlety at the cost of speed.

\section{Make failed matches and repeated work measurable}
\label{sec:ch85-measurement}
Rule-engine profiling should count more than successful rewrites.
Useful counters include candidate-index hits, complete matches, side-condition queries, rejected premises, stale records, committed edits, and requeued nodes.
A rule applied only once can still dominate compile time if its expensive premise is checked at every node.
Conversely, a frequently applied cheap identity may save downstream work by shrinking the graph early.
Measurements should retain the input corpus and pipeline configuration so these counts can be interpreted.

Reproducible scheduling helps diagnose regressions.
If two equally ranked matches are ordered by allocation address or hash-table iteration, unrelated changes can alter output and obscure the cause of a performance shift.
Assigning stable rule priorities and deterministic structural tie-breakers narrows that variability.
Parallel engines may deliberately permit nondeterministic schedules, but should then validate semantic invariants and report the reproducibility policy rather than implying byte-identical output.

Resource measurements also need representation awareness.
A persistent trace can retain old subgraphs after they disappear from the current root.
A node budget that counts only currently reachable nodes may therefore undercount actual memory.
The engine should either include retained history, caches, and private replacements in its accounting or state that the budget bounds only admitted IR size.
This is the same distinction that makes a product-state budget in the spellchecker different from a total-allocation limit.

\section{Shrink failures without changing the disputed premise}
\label{sec:ch85-shrink}
Randomly generated terms and graphs can expose interactions that hand-written patterns miss, but a large failing graph is difficult to diagnose.
A shrinking procedure removes unrelated substructure while preserving well-formedness and the observed disagreement.
For typed terms, replacing a subtree with a smaller term of the same type avoids manufacturing a new type error that masks the original bug.
For graph rewrites, shrinking must preserve the sharing, dominance, or use-count condition involved in the failure.
Turning a shared node into two copied trees can accidentally eliminate the very defect being investigated.

A useful failure record includes the original input, rule version, strategy, semantic mode, seed, budget, and independent oracle outcome.
The shrinker reruns the same oracle comparison after every proposed reduction of the test case.
If a smaller case fails only because it now exhausts fuel, that is not the same semantic counterexample unless budget exhaustion itself was the target bug.
Retaining this distinction prevents a fuzzing workflow from converting a concrete wrong result into an inconclusive timeout and calling the reduction successful.

The reference tests use exhaustive small families rather than a randomized shrinker.
Their tiny counterexamples are already inspectable, and fixed iteration order makes failures reproducible.
Adding randomized stress is a useful extension after the basic rule and certificate contracts are stable, not a substitute for establishing those contracts first.

\section{Exercises}
\label{sec:ch85-exercises}
\begin{enumerate}
\item\label{ex:ch85-priority} Two rules match the root but only the second rule's semantic premise holds.
Explain why root indexing cannot select the first rule unconditionally.
\item\label{ex:ch85-users} A shared add node has two parents; rewriting it changes a one-use predicate at both.
List the invalidation and worklist obligations.
\item\label{ex:ch85-stale} A queued match was found before a CFG edit.
Why must a cached dominance condition be revalidated or invalidated before commit?
\item\label{ex:ch85-budget} An expanding replacement exceeds the graph budget.
Specify a publication sequence that leaves the last committed IR valid.
\end{enumerate}

\section{Further reading}
\label{sec:ch85-reading}
Ohlebusch supplies the rewrite-theoretic context for strategy and normal-form claims~\cite{ohlebusch-2002-rewriting}.
The equality-saturation chapter and the \code{egg} paper describe an alternative-preserving engine with different update invariants~\cite{willsey-2021-egg}.
